%TODO
%Italizar
%revisar las referencias
%revisar todo cuando acabe la propuesta
%cambiar la imagen delas vistsas coordinadas
%actualizar consideraciones iniciales
%mas vis d tiempo
%interaccion
%explayarte en reducciones dimensionales
%Sacar vector caracteristicas
%LLMS escueto

\chapter{Marco Teórico}
\label{cap:marco_teorico}


\section{Consideraciones Iniciales}
Este capítulo presenta los conceptos teóricos que 
fundamentan la propuesta. Se comienza con \ac{VA}, 
su ciclo iterativo y el patrón de vistas coordinadas 
que organiza la interfaz. Luego, se revisa la 
visualización de procesos con dimensiones temporales 
y espaciales, y cómo la trazabilidad permite verificar 
los datos extraídos. Después se describen los 
\acs{LLM}, su uso para extraer información estructurada 
de texto, los problemas que esto implica y estrategias de mitigación. El capítulo 
cierra presentando los documentos institucionales 
fragmentados y el corpus 
utilizado para evaluar el sistema.

\section{\textit{Visual Analytics}}

\ac{VA} se define como la ciencia del razonamiento 
analítico facilitado por interfaces visuales 
interactivas~\cite{thomas2005illuminating}. Su núcleo 
es la integración del procesamiento automático con las 
capacidades de razonamiento humano, de modo que el 
analista participa activamente en el proceso de 
análisis y no solo observa sus resultados~\cite{keim2008visual, 
cui2019visual}.

A partir del modelo de referencia de Keim ~\cite{keim2008visual}, el proceso opera como un 
ciclo iterativo entre cuatro etapas interdependientes, 
tal como se ilustra en la Figura~\ref{fig:keim_va}. 
Primero, los datos crudos se transforman en estructuras 
manejables mediante tareas de limpieza y normalización. 
Luego toman dos vías simultáneas: la construcción de 
modelos automáticos que extraen patrones, y el mapeo de 
los resultados en representaciones visuales interactivas. 
A partir de esas representaciones, el analista evalúa 
los resultados, identifica patrones o anomalías y formula 
nuevas hipótesis. El 
conocimiento adquirido se reintegra al sistema mediante 
ajustes en los parámetros de los modelos o en los filtros 
de las vistas, reiniciando el ciclo.

\begin{figure}[H]
\centering
\includegraphics[width=\textwidth]{figs/CicloVA.pdf}
\caption{Ciclo de \ac{VA}, mostrando la interacción continua 
entre los datos, los modelos automáticos, las representaciones 
visuales y el analista. Adaptado de~\cite{keim2008visual}.}
\label{fig:keim_va}
\end{figure}

Sacha et al.~\cite{sacha2014} refinan este modelo 
distinguiendo explícitamente dos componentes que operan 
de forma complementaria. 
El componente computacional comprende los datos, los 
modelos analíticos y las visualizaciones generadas a 
partir de ellos. El componente humano comprende el 
proceso cognitivo del analista, que se estructura en 
tres bucles de razonamiento anidados.

El \textbf{bucle de exploración} (\textit{exploration loop}) 
describe cómo el analista interactúa con el sistema 
mediante acciones que producen hallazgos. Una acción 
puede ser filtrar un subconjunto de datos, cambiar el 
mapeo visual o ajustar los parámetros de un modelo. El 
sistema responde a esa acción y el analista observa el 
resultado, generando un hallazgo (\textit{finding}). Un 
hallazgo es cualquier observación relevante que el 
analista percibe en la visualización o en el modelo, 
como un patrón inesperado, un valor atípico o la 
ausencia de un dato esperado. Este bucle se repite de 
forma rápida y constituye la base de toda exploración 
interactiva.

El \textbf{bucle de verificación} (\textit{verification loop}) 
opera a un nivel superior. Cuando el analista interpreta 
un hallazgo en el contexto del dominio del problema, 
produce un \textit{insight}, una unidad de comprensión 
que conecta la evidencia visual con el conocimiento 
previo. Los \textit{insights} pueden confirmar o refutar 
una hipótesis existente, o generar hipótesis nuevas que 
requieran exploración adicional. Este bucle vincula lo 
que el analista observa en la interfaz con lo que sabe 
del dominio.

El \textbf{bucle de generación de conocimiento} 
(\textit{knowledge generation loop}) cierra el proceso. 
Cuando el analista acumula suficiente evidencia para 
confiar en un \textit{insight}, este se consolida como 
conocimiento. Ese conocimiento nuevo puede a su vez 
generar hipótesis adicionales, reiniciando el ciclo 
completo. Sacha et al. enfatizan que el conocimiento 
no se produce en el componente computacional ni en el 
humano por separado, sino en la interacción entre 
ambos~\cite{sacha2014}. El analista aporta juicio 
contextual y razonamiento que los algoritmos no poseen, 
mientras que los algoritmos procesan volúmenes de 
información que el analista no podría revisar 
manualmente. Esta complementariedad es la que 
permite que \ac{VA} funcione como un marco de generación 
de conocimiento.



La distinción entre estos tres bucles resulta relevante 
para el diseño de sistemas de \ac{VA} porque permite 
evaluar qué niveles de razonamiento soporta cada 
componente. Un sistema puede facilitar el 
bucle de exploración ofreciendo visualizaciones 
interactivas, pero si carece de mecanismos para que el 
analista verifique la evidencia, el bucle de verificación 
queda sin soporte. 





\subsection{\textit{Visual Analytics} aplicado a documentos 
textuales}
\label{sec:va_texto}

El procesamiento de información textual representa uno 
de los mayores desafíos dentro de \ac{VA} debido a la 
naturaleza no estructurada del lenguaje. 
Thomas y Cook  ~\cite{thomas2005illuminating} señalan que 
los documentos en texto libre requieren transformaciones 
semánticas profundas para extraer unidades de significado 
antes de poder ser representados visualmente o analizados 
algorítmicamente. A diferencia de datos numéricos o 
tabulares, donde las dimensiones del análisis están 
definidas de antemano, en un documento textual las 
entidades, las relaciones y la estructura temporal deben 
ser identificadas y extraídas como paso previo a 
cualquier representación visual. 
Kucher y Kerren~\cite{kucher2015text} proponen una 
taxonomía de técnicas de visualización de texto que 
organiza este campo según la tarea analítica, el tipo 
de dato y las propiedades de la representación visual. 
Esta taxonomía permite situar cualquier sistema de 
\ac{VA} para texto dentro de un espacio de diseño 
definido por las decisiones que conectan el dato textual 
con su forma visual.

Cuando los documentos provienen de dominios donde las 
decisiones tienen consecuencias institucionales, la 
extracción automática enfrenta una exigencia adicional. 
Los algoritmos que operan como cajas negras dificultan 
que el experto humano pueda auditar y confiar en los 
resultados. Hutchinson et al.~\cite{hutchinson2024llm} 
advierten que para incorporar modelos de lenguaje dentro 
de un entorno de \ac{VA}, sus respuestas deben estar 
ancladas a los datos originales. Esta necesidad de 
anclaje conecta directamente con el bucle de verificación 
de Sacha et al.~\cite{sacha2014}, donde el analista 
evalúa críticamente los hallazgos producidos por la 
máquina para determinar si constituyen evidencia real o 
un artefacto del procesamiento automático.

\subsubsection{Trazabilidad y lectura \textit{in situ}}
\label{sec:trazabilidad}

La trazabilidad es el mecanismo que vincula cada 
elemento visual de la interfaz con el fragmento 
específico del documento del que fue extraído. En 
términos del modelo de Sacha et al.~\cite{sacha2014}, 
este mecanismo soporta el bucle de verificación al 
permitir que el analista pase de un hallazgo visual a 
la evidencia textual que lo sustenta. Sin este vínculo, 
el analista solo puede aceptar o rechazar los resultados 
del modelo sin poder evaluar su corrección.

La lectura \textit{in situ} es la forma concreta en que 
se materializa la trazabilidad en la interfaz. Consiste 
en permitir que el analista, al seleccionar un dato en 
cualquiera de las vistas, acceda al fragmento exacto 
del documento original del que fue extraído, incluyendo 
su contexto circundante y su ubicación en la página. 
Este mecanismo permite un flujo de trabajo en el que 
el analista transita entre la vista panorámica de los 
datos y la lectura detallada del texto fuente sin 
abandonar la interfaz~\cite{thomas2005illuminating}.







\section{Visualización de Datos}
\label{sec:tecnicas_vis}

La extracción automática de información resulta 
insuficiente si el usuario no puede interpretar los 
resultados. La visualización de datos actúa como el 
canal que transforma información abstracta en 
representaciones gráficas, aprovechando la capacidad 
perceptiva humana para revelar patrones que de otro 
modo permanecerían ocultos en los 
datos~\cite{ortigossa2025time}. Esta sección presenta 
los principios de interacción, los patrones de 
coordinación entre vistas y las técnicas de 
representación que fundamentan el diseño de la 
interfaz.


\subsection{El Mantra de Shneiderman}
\label{sec:shneiderman}

El diseño de interfaces exploratorias se fundamenta en 
el principio propuesto por 
Shneiderman~\cite{shneiderman1996eyes}, conocido como 
el Mantra de Búsqueda de Información Visual. Este 
principio organiza la exploración en tres etapas 
jerárquicas. En la primera, \textit{overview}, el 
sistema presenta una perspectiva global del conjunto 
de datos que permite al analista orientarse sin 
saturación visual. En la segunda, \textit{zoom and 
filter}, el analista aísla subconjuntos de interés 
mediante controles interactivos, eliminando 
dinámicamente la información irrelevante. En la 
tercera, \textit{details-on-demand}, el analista 
selecciona un elemento específico y accede a su 
información completa. Este flujo descendente, de lo 
general a lo particular, permite que el analista 
construya progresivamente su comprensión del conjunto 
de datos sin necesidad de examinar cada elemento 
individualmente.


\subsection{Vistas múltiples coordinadas}
\label{sec:mcv}

Cuando los datos poseen múltiples dimensiones, una sola 
representación visual rara vez es suficiente para 
explorarlos por completo. Las vistas múltiples coordinadas (\acs{MCV}, pos sus siglas en ingles) son un patrón de 
diseño en el que dos o más representaciones visuales del 
mismo conjunto de datos se presentan de forma simultánea, 
y las acciones del usuario en una vista se propagan 
automáticamente a las 
demás~\cite{roberts2007state, baldonado2000guidelines}. 
De este modo, cada vista puede especializarse en una 
dimensión del análisis mientras la coordinación entre 
ellas preserva el contexto compartido 
(Figura~\ref{fig:brushing_linking}).

El mecanismo fundamental de coordinación es la selección y vinculación interactiva 
(\textit{brushing and linking})~\cite{becker1987brushing}. 
Cuando el analista selecciona un elemento en una vista, 
esa selección se transmite a todas las vistas vinculadas, 
donde el mismo elemento queda resaltado en su 
representación correspondiente. Un ejemplo de esto puede observarse en la 
(Figura~\ref{fig:brushing_linking}). 
Becker y Cleveland~\cite{becker1987brushing} definen 
cuatro operaciones de \textit{brushing} (resaltado, 
resaltado con sombra, eliminación y etiquetado) cuyo 
efecto se propaga simultáneamente a todas las vistas. 
Este mecanismo permite que el analista observe cómo se 
manifiesta una misma entidad en distintas dimensiones 
sin perder la referencia de qué elemento está analizando.

% FIGURA: Adaptar de Becker & Cleveland 1987, Figure 1
% Muestra scatterplots con brushing: seleccionar puntos 
% en un scatterplot los resalta en los demás.
% También puedes adaptar de Roberts 2007, Figure 1 
% (Improvise) que muestra 6 vistas coordinadas.
\begin{figure}[H]
\centering
\includegraphics[width=0.85\textwidth]%
{figs/improvise.png}
\caption{Sistema Improvise mostrando seis vistas 
coordinadas del mismo conjunto de datos. La selección 
de elementos en una vista se propaga a las demás 
mediante \textit{brushing and linking}. Imagen 
extraída de~\cite{roberts2007state}.}
\label{fig:brushing_linking}
\end{figure}

Además del \textit{brushing}, 
Roberts~\cite{roberts2007state} identifica otros 
mecanismos de coordinación entre vistas. La navegación 
sincronizada permite que un cambio de zoom o 
desplazamiento en una vista se replique en las demás. 
El filtrado compartido permite que los criterios de 
exclusión aplicados en una vista reduzcan el conjunto 
visible en todas las vistas vinculadas. La combinación 
de estos mecanismos permite al analista explorar 
distintas facetas de los datos de forma fluida, 
manteniendo la coherencia entre las representaciones.

Desde una perspectiva computacional, la coordinación 
entre vistas implica mantener un estado compartido de 
selección y filtrado que se actualiza y propaga cada 
vez que el usuario interactúa con cualquiera de las 
vistas. 
%North y Shneiderman~\cite{north2000snap}  formalizan este mecanismo utilizando el modelo relacional de datos, donde las relaciones entre registros de distintas tablas determinan qué elementos se vinculan entre vistas. Esta formalización permite definir la coordinación de forma independiente de las representaciones visuales específicas que se utilicen.

Baldonado et al.~\cite{baldonado2000guidelines} 
proponen directrices para el diseño de estos sistemas. 
Entre ellas destacan el principio de diversidad, 
que establece que las vistas deben usar representaciones 
distintas para aportar perspectivas complementarias 
sobre los mismos datos, y el principio de 
parsimonia, que recomienda usar el menor número 
de vistas necesario para cubrir las dimensiones del 
análisis sin imponer una carga cognitiva excesiva al 
usuario. Estos principios implican que cada vista 
adicional debe justificarse por la perspectiva nueva 
que aporta, y que la coordinación entre ellas debe 
reducir, y no aumentar, el esfuerzo cognitivo del 
analista.

En el marco de \ac{VA}, las \acs{MCV} cumplen un rol 
especialmente importante porque permiten sostener los 
bucles de razonamiento descritos por 
Sacha et al.~\cite{sacha2014} a través de múltiples 
dimensiones de forma simultánea. El bucle de 
exploración se beneficia porque el analista puede 
generar hallazgos comparando lo que observa en una 
vista con lo que aparece en otra. El bucle de 
verificación se beneficia porque un dato sospechoso 
en una vista puede contrastarse inmediatamente con 
su representación en una vista diferente.

\subsection{Representación de datos temporales}
\label{sec:temporal}

El tiempo posee una estructura semántica que lo 
distingue de otras variables cuantitativas. 
Aigner et al.~\cite{aigner2011visualization} proponen 
una taxonomía que caracteriza los datos temporales 
según tres propiedades. La escala define si el tiempo 
se representa de forma ordinal (solo el orden importa) 
o cuantitativa (las distancias entre instantes son 
medibles). El alcance distingue entre datos puntuales, 
que registran un instante, y datos basados en 
intervalos, que poseen inicio y fin. El arreglo 
describe si la progresión temporal es lineal o cíclica. 
Estas propiedades determinan qué técnicas de 
representación son adecuadas para cada tipo de dato 
temporal. Ortigossa et al.~\cite{ortigossa2025time} 
ofrecen una revisión amplia de técnicas de 
visualización para series temporales, organizadas 
según estas propiedades y según el tipo de tarea 
analítica que soportan.

\subsubsection{Gráficos de líneas}
\label{sec:linechart}

El gráfico de líneas es la forma más antigua de 
representación de series temporales. 
Playfair~\cite{playfair1786atlas} introdujo esta 
técnica en 1786 para representar la evolución de las 
exportaciones e importaciones de Inglaterra a lo largo 
del siglo XVIII, mapeando el tiempo al eje horizontal 
y los valores al eje vertical. Dos siglos después, 
esta representación permanece esencialmente sin 
cambios~\cite{tufte2001visual}. Cuando se superponen 
múltiples series en el mismo espacio, el analista puede 
comparar tendencias, identificar puntos de cruce y 
detectar divergencias en la evolución de distintas 
variables. Cleveland~\cite{cleveland1993visualizing} 
establece principios perceptivos para la lectura 
efectiva de estos gráficos, como la elección de la 
relación de aspecto (\textit{aspect ratio}) para 
facilitar la percepción de las pendientes. La principal 
limitación del gráfico de líneas aparece cuando el 
número de series crece, ya que la superposición 
dificulta la lectura individual de cada una. La 
Figura~\ref{fig:line_chart} ilustra la superposición 
de tres series temporales sobre un eje común.

\begin{figure}[H]
\centering
\includegraphics[width=0.8\textwidth]{figs/line_chart.png}
\caption{Gráfico de líneas mostrando la evolución 
temporal de tres variables a lo largo de doce meses. 
La superposición de series permite comparar tendencias 
y detectar divergencias.}
\label{fig:line_chart}
\end{figure}

\subsubsection{Diagramas de Gantt}
\label{sec:gantt}

Para datos basados en intervalos con arreglo lineal, 
los diagramas de Gantt constituyen la representación 
visual más adecuada~\cite{aigner2011visualization, 
ortigossa2025time}. Desarrollado originalmente por 
Henry Gantt a inicios del siglo 
XX~\cite{wilson2003gantt}, este diagrama asocia cada 
entidad a una fila horizontal y representa los 
intervalos mediante barras cuya longitud es proporcional 
a su duración. Esta técnica permite comparar 
visualmente solapamientos entre fases, identificar 
retrasos y evaluar la progresión relativa de múltiples 
procesos sobre un eje temporal 
común~\cite{aigner2011visualization}. La 
Figura~\ref{fig:gantt_time} muestra cinco entidades 
con intervalos de distinta duración y solapamiento.

\begin{figure}[H]
\centering
\includegraphics[width=0.85\textwidth]{figs/gantt_chart.png}
\caption{Diagrama de Gantt mostrando los intervalos 
temporales de cinco entidades. La longitud de cada 
barra es proporcional a la duración del intervalo, 
y la disposición en filas permite identificar 
solapamientos y retrasos entre procesos paralelos. }
\label{fig:gantt_time}
\end{figure}

\subsubsection{\textit{Streamgraph}}
\label{sec:streamgraph}

El \textit{streamgraph} es una variante del gráfico de 
áreas apiladas en la que múltiples series temporales se 
representan como flujos cuyo grosor varía en función del 
valor de cada variable en cada punto del tiempo. 
Havre et al.~\cite{havre2002themeriver} introdujeron 
esta técnica con el nombre de ThemeRiver para 
visualizar la evolución temática de grandes colecciones 
de documentos. Cada tema se representa como una 
corriente cuyo ancho refleja su frecuencia, y el 
conjunto de corrientes forma un flujo continuo a lo 
largo del eje temporal. Byron y 
Wattenberg~\cite{byron2008stacked} formalizaron la 
geometría del \textit{streamgraph}, proponiendo 
algoritmos para optimizar la disposición de las capas 
y minimizar la distorsión visual que surge al apilar 
múltiples series.

A diferencia del gráfico de líneas, que enfatiza los 
valores individuales de cada serie, el 
\textit{streamgraph} enfatiza las proporciones 
relativas y la composición del total a lo largo del 
tiempo~\cite{byron2008stacked}. Esto lo hace 
especialmente útil cuando el analista necesita percibir 
cómo cambia la distribución de un conjunto de variables 
entre períodos. La Figura~\ref{fig:streamgraph} muestra 
cómo cuatro variables compiten en proporción a lo 
largo del tiempo.

\begin{figure}[H]
\centering
\includegraphics[width=0.85\textwidth]{figs/streamgraph.png}
\caption{\textit{Streamgraph} mostrando la evolución 
de cuatro variables a lo largo del tiempo. El grosor 
de cada flujo refleja el valor de la variable en cada 
punto, y la disposición simétrica facilita la 
percepción de las proporciones relativas entre 
variables.}
\label{fig:streamgraph}
\end{figure}

\subsubsection{Diagramas de Sankey}
\label{sec:sankey}

Mientras que el \textit{streamgraph} muestra cómo 
cambia la composición de un conjunto a lo largo del 
tiempo, el diagrama de Sankey muestra cómo fluyen 
cantidades entre estados o categorías. Originalmente 
concebido por el capitán irlandés Matthew 
Sankey en 1898 para representar las pérdidas 
energéticas de una máquina de 
vapor~\cite{schmidt2008sankey}, este tipo de diagrama 
utiliza nodos para representar estados y conexiones 
cuyo ancho es proporcional a la cantidad que transita 
entre ellos. Riehmann et al.~\cite{riehmann2005sankey} 
proponen un sistema interactivo que soporta el flujo 
de exploración de Shneiderman, ofreciendo una vista 
general del grafo de flujos con la posibilidad de 
hacer zoom y acceder a detalles bajo demanda.

Cuando se incorpora la dimensión temporal, los nodos 
se organizan en columnas que representan momentos 
sucesivos, y los flujos entre columnas muestran las 
transiciones entre estados de un momento al siguiente. 
Esta disposición permite identificar qué transiciones 
son frecuentes, dónde se concentran las pérdidas y 
qué trayectorias siguen las entidades a lo largo del 
proceso. La Figura~\ref{fig:sankey} ilustra este 
esquema con tres estados distribuidos en tres momentos 
sucesivos.

\begin{figure}[H]
\centering
\includegraphics[width=0.85\textwidth]{Tesis_CS_202203/figs/sankey_horizontal.png}
\caption{Diagrama de Sankey temporal con tres estados 
(A, B, C) en tres momentos sucesivos ($t_1$, $t_2$, 
$t_3$). El ancho de cada flujo es proporcional a la 
cantidad de entidades que transitan entre estados, 
permitiendo identificar las trayectorias más 
frecuentes.}
\label{fig:sankey}
\end{figure}

\subsection{Representación geoespacial}
\label{sec:geoespacial}

Cuando los datos poseen una dimensión geográfica, las 
técnicas de geovisualización permiten revelar patrones 
espaciales que permanecen invisibles en representaciones 
tabulares. Andrienko y Andrienko~\cite{andrienko2006exploratory} 
establecen que el análisis exploratorio de datos 
espaciales requiere técnicas capaces de representar 
simultáneamente la distribución geográfica de los datos 
y sus atributos temáticos, de modo que el analista pueda 
percibir tanto dónde ocurren los fenómenos como cuál es 
su magnitud.

Las dos técnicas fundamentales de la cartografía temática 
son los mapas coropléticos y los mapas de símbolos 
proporcionales~\cite{slocum2014thematic}. Los mapas 
coropléticos codifican valores agregados mediante la 
variación de color o intensidad en regiones predefinidas, 
lo que los hace adecuados para representar densidades, 
tasas o proporciones cuya distribución sigue límites 
administrativos. Los mapas de símbolos proporcionales 
representan magnitudes mediante marcadores cuyo tamaño 
varía según el valor asociado a cada ubicación, lo que 
los hace preferibles cuando los datos corresponden a 
conteos o totales asociados a puntos específicos. Ambas 
técnicas pueden combinarse en un mismo mapa: las regiones 
aportan el contexto geográfico mediante color, mientras 
que los símbolos superpuestos comunican magnitudes 
puntuales~\cite{slocum2014thematic}. A este tipo de 
representación, que integra polígonos coloreados, 
símbolos y etiquetas en una sola vista, se le denomina 
mapa anotado.

Cuando las entidades se distribuyen en múltiples niveles 
administrativos, la agregación jerárquica permite 
transitar entre escalas. En el nivel más amplio el mapa 
muestra valores agregados por región. Al hacer zoom, los 
agregados se descomponen en las unidades que los conforman, 
revelando la distribución interna. Este mecanismo, conocido 
como \textit{zoom semántico}, ajusta el nivel de detalle 
de la representación según la escala de 
visualización~\cite{andrienko2006exploratory}. La 
Figura~\ref{fig:geo_zoom} ilustra este esquema en tres 
niveles: nacional, regional y local.

% FIGURA: Tres paneles horizontales mostrando zoom semántico
% Panel 1 (nacional): mapa con regiones coloreadas
% Panel 2 (regional): zoom a una región con símbolos proporcionales
% Panel 3 (local): zoom a una ubicación con anotaciones de detalle
\begin{figure}[H]
\centering
\includegraphics[width=\textwidth]{figs/geo_zoom.png}
\caption{Zoom semántico en tres niveles sobre un mapa 
anotado. En el nivel nacional las regiones se colorean 
según un valor agregado. Al hacer zoom a una región 
aparecen símbolos proporcionales por entidad. En el 
nivel local se muestran las anotaciones de detalle de 
una entidad específica.}
\label{fig:geo_zoom}
\end{figure}

La combinación del \textit{zoom semántico} con los 
principios de Shneiderman permite implementar el flujo de \textit{overview}, 
\textit{zoom and filter} y \textit{details-on-demand} 
en la dimensión espacial. El nivel nacional corresponde 
al \textit{overview}; la selección de una región activa 
el \textit{zoom and filter}; y la selección de una 
entidad específica despliega sus \textit{details-on-demand}. 
Esta correspondencia entre los principios de interacción 
y los niveles geográficos convierte al mapa anotado en 
un punto de entrada natural para la exploración de 
datos distribuidos geográficamente.

\subsection{Reducción dimensional y proyección semántica}
\label{sec:reduccion}

Cuando cada entidad de un conjunto de datos se 
describe mediante un vector de alta dimensión, la 
comparación directa entre entidades resulta 
cognitivamente inviable. Las técnicas de reducción 
dimensional proyectan estos vectores en un espacio de 
dos o tres dimensiones preservando, en la medida de lo 
posible, las relaciones de proximidad del espacio 
original. En la proyección resultante, las entidades 
que son similares en el espacio de alta dimensión 
aparecen cercanas en el plano, permitiendo al analista 
identificar agrupaciones, valores atípicos y patrones 
de similitud de forma visual.

t-SNE (\textit{t-distributed Stochastic Neighbor 
Embedding})~\cite{vandermaaten2008tsne} fue una de las 
primeras técnicas en lograr proyecciones de alta calidad 
para visualización, modelando las similitudes entre 
puntos como distribuciones de probabilidad y 
minimizando la divergencia entre las distribuciones 
del espacio original y el proyectado. Sin embargo, 
t-SNE presenta limitaciones en la preservación de la 
estructura global de los datos y en su escalabilidad 
a conjuntos de datos grandes.

\acs{UMAP} (\textit{Uniform Manifold Approximation and 
Projection})~\cite{mcinnes2018umap} aborda estas 
limitaciones. Fundamentado en geometría riemanniana 
y topología algebraica, \acs{UMAP} construye una 
representación topológica difusa del espacio de alta 
dimensión y optimiza su proyección en dimensiones 
bajas minimizando la entropía cruzada entre ambas 
representaciones. Comparado con t-SNE, \acs{UMAP} preserva 
mejor la estructura global de los datos, escala a 
conjuntos de datos significativamente más grandes y 
no tiene restricciones computacionales sobre la 
dimensión del espacio 
proyectado~\cite{mcinnes2018umap}.

Para que la reducción dimensional sea aplicable a 
datos textuales, los documentos o fragmentos de texto 
deben primero transformarse en vectores numéricos de 
alta dimensión. Los modelos de 
\textit{embeddings} textuales, como Sentence 
Transformers~\cite{reimers2019sbert}, proyectan 
oraciones o párrafos a un espacio vectorial denso en 
el que la distancia entre vectores refleja la 
similitud semántica entre los textos que representan. 
La combinación de \textit{embeddings} textuales con 
reducción dimensional mediante \acs{UMAP} permite visualizar 
grandes conjuntos de documentos como nubes de puntos 
en las que la proximidad espacial corresponde a la 
similitud de contenido.




\section{Procesamiento de Lenguaje Natural con \texorpdfstring{\acs{LLM}}{LLMs}}
La extracción de datos en documentos legales requería programar reglas manuales o entrenar modelos con miles de ejemplos específicos \cite{zhong-etal-2020-nlp}. Esta técnica resulta inviable para procesar expedientes de control gubernamental, cuyas estructuras y narrativas varían constantemente. Frente a este obstáculo, la presente investigación adopta los Grandes Modelos de Lenguaje como motor de extracción, aprovechando su capacidad para procesar semántica compleja sin requerir un reentrenamiento masivo por cada nuevo tipo de documento.


\subsection{Modelos y Extracción de Información}
Los \acs{LLM} poseen la capacidad de realizar tareas de extracción semántica avanzada sin requerir un reentrenamiento masivo. Mediante la técnica de aprendizaje con pocos ejemplos (\textit{few-shot learning})\cite{brown2020}. 

Para asegurar que la extracción de hitos, fechas y presupuestos mantenga una secuencia lógica, se emplea la estrategia de \textit{Prompting} denominada Cadena de Pensamiento (\textit{Chain-of-Thought}). Esta técnica induce al \acs{LLM} a descomponer un problema de análisis complejo en pasos de razonamiento intermedios antes de emitir una respuesta final, mejorando significativamente la fidelidad de la información extraída a partir de textos densos \cite{wei2022}.

\begin{figure}[H]
\centering
\includegraphics[width=0.8\textwidth]{Tesis_CS_202203/chain_of_thought.png}
\caption{Comparación entre el prompting estándar y el razonamiento por Cadena de Pensamiento, el cual muestra con claridad los pasos intermedios. Imagen extraída de \cite{wei2022}.}
\label{fig:cot_prompting}
\end{figure}



\subsection{Alucinaciones y Trazabilidad}
A pesar de su capacidad de razonamiento, la implementación de \acs{LLM} en entornos de alta responsabilidad enfrenta un riesgo crítico: las alucinaciones. \cite{ji2023} definen este fenómeno como la generación de contenido irreal que aparenta ser verdadero, clasificándolo en dos tipos: alucinaciones intrínsecas (cuando el texto generado contradice el documento fuente) y extrínsecas (cuando se genera información que no puede ser verificada en la fuente original). En una auditoría, alucinar la fecha , presupuesto o el lugar invalida por completo el análisis.

Frente al riesgo de las alucinaciones, la extracción de información no puede operar como una caja negra. \cite{hutchinson2024llm} advierten que para confiar en un \acs{LLM} dentro de un entorno de \textit{Visual Analytics}, sus respuestas deben estar ancladas a los datos originales. En el trabajo diario de un auditor, esto se traduce en una necesidad: el sistema debe crear un enlace visual directo entre el dato extraído por la máquina y el párrafo exacto del expediente. De esta forma, el analista no asume automáticamente que el modelo tiene la razón, sino que utiliza la interfaz para auditar la fuente, corregir errores y validar que la información sea real.

\subsection{Estructuración de Narrativas Complejas}
El objetivo de extraer entidades, tiempos y atributos no es generar bases de datos aisladas, sino reconstruir una narrativa compleja. En el contexto de esta investigación, una narrativa compleja se define como la progresión cronológica y causal de eventos burocráticos distribuidos a través de múltiples documentos (informes, adendas, resoluciones). 

Para modelar computacionalmente esta narrativa, el sistema extrae eventos. Un evento se define como una tupla estructurada que responde al qué, quién y cuándo de una acción documentada \cite{navas2022whenthefact, yao2022leven}. Al transformar el texto fragmentado en una secuencia de eventos, el \textit{Visual Analytics} logra mapear el flujo de la información. De esta forma, la interfaz revela visualmente cómo evolucionan los procesos y dónde se concentran las demoras, eliminando la necesidad de leer miles de documentos de forma secuencial.


\section{Consideraciones Finales}
El marco teórico desarrollado en este capítulo establece los pilares conceptuales para abordar el análisis de documentos gubernamentales. En primer lugar, se ha definido que el \textit{Visual Analytics} provee la arquitectura metodológica necesaria para mantener al analista humano en el centro del proceso de verificación. Paralelamente, las técnicas de visualización espaciotemporal y multidimensional proporcionan los mecanismos para representar la estructura cronológica y las métricas de las auditorías sin generar saturación visual.

Por su parte, se estableció que los Grandes Modelos de Lenguaje actúan como el motor de extracción capaz de procesar la semántica del texto y estructurar las narrativas complejas. Sin embargo, su implementación en escenarios de alta responsabilidad exige mecanismos de trazabilidad visual para mitigar el riesgo de alucinaciones y garantizar la fidelidad de la información.

En el siguiente capítulo se presentara el Estado del Arte. Se revisarán los trabajos enfocados en sistemas de extracción legal, herramientas visuales para texto y aplicaciones de los \acs{LLM}..
